\section{Hạn chế}

Trong phần Hạn chế của nghiên cứu, cần làm rõ những yếu tố có thể ảnh hưởng đến độ đầy đủ, tính chính xác và khả năng khái quát hóa của kết quả. Các hạn chế có thể bao gồm phạm vi và nguồn dữ liệu, khi nghiên cứu chủ yếu dựa trên dữ liệu từ Retraction Watch Database kết hợp với Scopus và Web of Science; tính không đồng nhất hoặc thiếu hụt của một số trường thông tin, đặc biệt là lý do và thời điểm thu hồi; ảnh hưởng của các sự kiện thu hồi hàng loạt đến phân bố số lượng ấn phẩm theo năm và các chỉ số trắc lượng; cũng như độ trễ giữa thời điểm công bố và thời điểm thu hồi, khiến các ấn phẩm được công bố trong những năm gần đây có thể chưa được quan sát đầy đủ. Bên cạnh đó, các chỉ số trắc lượng và phân tích mô tả chủ yếu phản ánh các đặc điểm và mối liên hệ trong dữ liệu, do đó không đủ cơ sở để xác lập quan hệ nhân quả giữa các yếu tố liên quan đến việc thu hồi. Vì vậy, các kết quả cần được diễn giải trong phạm vi bộ dữ liệu, nguồn dữ liệu và giai đoạn 2021 - 2026 được nghiên cứu, thay vì khái quát trực tiếp cho toàn bộ hiện tượng thu hồi ấn phẩm trong lĩnh vực Khoa học thông tin.